\section{Source reconciliation and further research}
\label{audit:section}

\subsection{Corrections, qualifications, and status changes}

\paragraph{Current source status.}
The targeted audit did not establish an additional false theorem in the
current source dependencies. The useful corrections are therefore
precise integration and scope corrections, not a claim to have found
a new error merely because a stronger statement is unavailable.
The two Stieltjes corrections already incorporated in the unified
manuscript must be retained. Its chapter
\path{11-stieltjes-source-corrections.tex} explicitly repairs an
initial-index discrepancy in the inspected version of a harmonic-zeta
paper and a literal-coefficient sign description in a height-one
multiple-zeta paper \cite{ProveIt}. Neither is a new discovery of this
delivery.

The exact conventions can be checked without a numerical fit. From
\eqref{intro:stieltjes} and the Hurwitz recurrence,
\[
 \gamma_r(a+1)=\gamma_r(a)-\frac{\log^r a}{a}.
\]
Consequently a Stieltjes limit whose sum begins at $n=1$ instead of
$n=0$ has lost $\log^r(a)/a$. For $r=0,a=1$ the lost quantity is
one. Likewise the literal coefficient of $u$ in
$\zeta(1+u,a)$ is $-\gamma_1(a)$, rather than $\gamma_1(a)$.
These facts explain both inherited corrections and fix the sign
convention needed for \eqref{ats:eq:stieltjes} and
\eqref{sd:eq:resonant}.

\paragraph{Higher-order pole terms.}
For a product $C_k(s)\zeta(1+2s,a)$, the coefficient of $s^r$
uses the coefficient of $s^{r+1}$ in $C_k$. Formula
\eqref{sd:eq:resonant} retains this term as
$r!\alpha_{k,r+1}/2$. A proposed extension obtained by differentiating
only the regular Laurent part would be incorrect. The present article
supplies the complete replacement, without suggesting that the
source's proved first-derivative theorem was wrong.

\paragraph{Ordinary sums and continued values.}
The continuous-order square law is an ordinary identity only on its
stated half-plane. At $s=0$ and a fixed positive $a$, continuation gives
$\mathcal B_0(z)=1/2$. The continued lower-end expression in
\eqref{ats:eq:continuous} is then $1/4$, although every pointwise
summand of its left side at $s=0$ is zero. In the finite identity
\eqref{ats:eq:continuous-finite}, the upper endpoint is exactly $1/4$
and cancels it. This is a concrete obstruction to silently extending
the ordinary summation prescription across that boundary. Similarly,
the finitely many low-order coefficients in
\eqref{mm:eq:confluence} are continued spectral values when the
unmodified weighted series diverges. The article states that distinction
explicitly.

\paragraph{Moment reconstruction.}
The gap-four theorem is already proved in \cite{HTGC,HTST}. An older
research-status paragraph that still lists it as wholly open should
be updated by citing those reports. For order eight, the inherited
theorem outside $b\equiv1\pmod{16}$ is strengthened here; it is not
relabelled as an original result. The new statements are the smaller
exceptional class in Theorem~\ref{rt:rank8} and the order-ten family in
Theorem~\ref{rt:rank10}.

Experimental valuation laws must be replaced by their proved
finite-precision versions. For example, the new order-eight formula
gives valuations $4$ at $b=65$ and $2$ at $b=129$, while the
order-ten formula gives valuation $2$ at $b=66$.
These follow from $b-8001$ and $b-962$, respectively. They do not
justify a global law centered at an ordinary integer: the excluded
progressions remain outside the theorem. Moreover, formal full rank
refers to the specified moment/zeta coordinate module. It must not
be converted into arithmetic independence of its numerical evaluations.

\paragraph{Cyclotomic parity.}
The finite coefficient identity in Section~\ref{cyc:sec} is a form
of the classical colored depth-two parity theorem \cite{PanzerParity}.
The addition is its explicit rational-tail resummation with arbitrary
outer character, its cutoff proof, and the specialized evaluations.
The complementary conjugation sector is not reduced. Conditional
bilateral intermediate sums are locally uniformly convergent; the
projected double sums used in the final generator are absolutely
normally convergent. Replacing the first statement by absolute normal
convergence would be an unjustified strengthening.

\paragraph{Gaussian conjectures and formal obstructions.}
The source's Gaussian weight-six and revised weight-eight conjectures
remain unresolved here. A non-membership result in a specified
shuffle, parity, or Cayley relation space is a limitation of that
relation system. It is not a disproof of an equality between actual
periods. The newly evaluated Gaussian rational tails have their own
definitions and cannot be substituted for those older conjectural
quantities. No integration note should mark the older problems solved
on the strength of the present conductor-four examples.

\subsection{A proposed research programme with concrete endpoints}

Each question below asks for an identity, a complete structural
description, or an exact nonvanishing certificate. The order reflects
dependence on the proved results, not a claim about difficulty.

\begin{question}[The two remaining moment progressions]
Can one prove nonvanishing in the excluded classes
$b\equiv8001\pmod{16384}$ for order eight and
$b\equiv962\pmod{2048}$ for order ten?
Lemma~\ref{rt:general-B} makes higher congruence precision effective.
A useful next result would either control the exceptional classes
uniformly at every precision or give an independent exact argument on
them. Certifying any fixed number of additional binary digits still
leaves an infinite progression. One must also distinguish a compatible
$2$-adic root from an ordinary positive integer at which the rational
determinant actually vanishes.
\end{question}

\begin{question}[The full gap-six family]
For every even $p=2a$ and $q=2a+6$, does the first-$q$ prefix have
full formal rank? The finite example $(p,q)=(10,16)$ is included here,
but a proof uniform in $a$ is still missing. The fixed-order contraction
has the threshold $b\ge2a$, whereas here $b=a+3$. For $a>3$ its
omitted $B_0$ and $B_1$ contributions cannot be discarded. An exact
bordered-Hankel formula that retains those moving endpoint terms is
a concrete next target.
\end{question}

\begin{question}[All unequal even orders]
Find a uniform nonvanishing theorem for $M_{a,b}$, $1\le a<b$,
or identify and prove the exact exceptional set if one exists.
A useful intermediate result would characterize the small determinant
after the Pad\'e transformation in terms of explicit orthogonal
polynomial evaluations and their finite differences. The endpoint
should be an exact identity or sign/nonzero certificate, rather than
an extrapolation from a large table of ranks.
\end{question}

\begin{question}[Bilinear adjacent-tail certificates]
Replace $E_n^2$ in Theorem~\ref{ats:thm:polynomial} by a product of
tails with distinct polynomial denominators $P,Q$. Allow both products
of adjacent tails in the certificate. Which rational coefficient
functions eliminate every linear-tail term, and which pairs $P,Q$
admit a finite endpoint evaluation? The compatibility calculation
in Section~\ref{ats:section} classifies only the displayed one-product
ansatz. The first nontrivial test is two distinct quadratic
resolvents with independent shifts.
\end{question}

\begin{question}[Continuation of the ordinary square law]
Continue \eqref{ats:eq:continuous} across $\Re s=0$ with a fully
specified subtraction of its exact upper endpoint. Determine the
contact term and all order derivatives at each nonpositive integer.
The $s=0$ discrepancy of $1/4$ described above is a mandatory first
test. An all-index formula would connect negative-order alternating
Hurwitz values, Euler polynomials, and Laurent coefficients without
confusing ordinary sums with continued values.
\end{question}

\begin{question}[Independent shifts and independent order jets]
The two factors of $P(k)=(k+a)(k+b)$ in
\eqref{ats:eq:two-shift} belong to one tail. They do not answer the
source question about
$\zeta(p,n+v)\zeta(q,n+w)$ with independently shifted tails.
Derive a finite-cutoff identity for those products, then differentiate
independently in $p,q$, and the outer weight order. The desired formula
must state all moving polar hyperplanes and their residues. This would
extend the matched continuous-order diagonal treated here and resolve
the corresponding independent-shift question in \cite{HTST}.
\end{question}

\begin{question}[The complementary cyclotomic sector]
Construct the two-resolvent generator for
\[
 \frac{R_{a,b;c}(x,y)-R_{a^{-1},b^{-1};c^{-1}}(x,y)}{2i}
\]
in the odd-weight setting. A first conductor-four coefficient is
$\Im Z(2,1;i,-i)$. Specify a candidate collection of classical or
cyclotomic polylogarithms and an exact reduction into it, or retain a
precisely defined additional coordinate. The parity generator alone
does not determine this sector. Failure of a finite numerical search
would not prove independence from the proposed collection.
\end{question}

\begin{question}[Shifted and degree-generating cyclotomic normal forms]
Introduce independent rational shifts into the bilateral lattice proof
and compute its finite boundary terms. Then introduce a third generating
variable that sums the outer polynomial degree $h$ in
\eqref{cyc:allweights}. Having a formula for every individual $h$ does
not yet produce that analytic generating function. A compact answer
would organize all resolvent and weight derivatives by a finite
Bell-polynomial rule, with its convergence region stated.
\end{question}

\begin{question}[Three-tail products]
Develop a three-tail analogue of the finite outer-cutoff decomposition.
The natural target is an explicit depth-two remainder plus finite
single-polylogarithm terms and endpoints. Determine the exact parity
sector in which this reduction occurs. The product of three tails
has a new ordered region structure, so unrestricted depth-one closure
would be an unsupported extrapolation from Section~\ref{cyc:sec}.
\end{question}

\begin{question}[A finite description of the logarithmic Gamma correction]
At which rational shifts and nonzero masses can $S(a,b)$ in
\eqref{gp:eq:S} be expressed in a stated finite collection of
polylogarithm values and order derivatives? The elementary-kernel
forms \eqref{gp:eq:integer}--\eqref{gp:eq:half} and the exact
shift/distribution rules provide starting equations. A rigorous
finite evaluation at a single nontrivial mass would be progress.
For generic parameters, the Gamma integral remains a complete exact
representation; no minimality theorem for that representation has
been established.
\end{question}

\begin{question}[Several masses at every spectral order]
Extend Theorem~\ref{mm:thm:second} to $r\ge3$. Pointwise, the variance
identity is replaced by a higher symmetric polynomial in the logarithms
$L_j$. Determine the finite-subtracted kernels, their confluent
coefficients, and their mixed mass primitives. With even weights
there is no pole at zero; with odd weights, compute the additional
Stieltjes and one-higher-coefficient residue terms rather than assuming
the even-weight prescription still applies.
\end{question}

\begin{question}[Shift, mass, and order derivatives in one formula]
Combine the primitive tower $J_m$ with the divided-difference kernel
$R_{P,a}$. Can every rectangular mixed shift--mass primitive of the
second defect be written using a finite, explicitly anchored collection
of Hurwitz derivatives and ordinary integrals? The mixed derivatives
must commute on a stated domain, and constants must be fixed at all
base faces. This is a stronger target than the separate one-variable
and double-mass primitives proved here.
\end{question}

\begin{question}[Joint tail identities and spectral identities]
The square certificate, the cyclotomic parity generator, and the
quadratic resolvent all produce polygamma expressions. At a fixed
conductor and weight, form their joint exact relation system and
determine which mixed harmonic-tail coefficients it isolates. A
particularly concrete endpoint is a new closed formula for a single
mixed tail product, rather than only the convolution in
\eqref{ats:eq:convolution}. Keep the formal relation quotient distinct
from the space of numerical periods.
\end{question}

\begin{question}[Small trusted certificates and formalization]
Which of the finite proofs can be reduced to short rational identities
that a proof assistant can check directly? The adjacent-tail recurrence,
the harmonic coefficient formula, the resonant Laurent multiplication,
and the two modular substitution certificates are natural first
targets. The analytic continuation lemmas then need explicit locally
uniform majorants and an agreed special-function library. A formal
check of the finite data alone would not formalize the associated
infinite-dimensional theorem.
\end{question}

